\documentclass[11pt]{article}
\usepackage[a4paper,margin=1in]{geometry}
\usepackage{amsmath,amssymb,amsthm,mathtools}

\newtheorem{theorem}{Theorem}[section]
\newtheorem{lemma}[theorem]{Lemma}
\newtheorem{proposition}[theorem]{Proposition}
\newtheorem{conjecture}[theorem]{Conjecture}
\newtheorem{remark}[theorem]{Remark}

\newcommand{\Om}{\Omega}
\newcommand{\R}{\mathbb R}
\newcommand{\T}{\mathcal T}
\newcommand{\A}{\mathcal A}
\newcommand{\dT}{\delta_T}
\newcommand{\HH}{\mathcal H}

\title{Hybrid Level-Flow Graph Entry: Repair Attempt and Exact Remaining Gap}
\author{}
\date{}

\begin{document}
\maketitle

\section{The intended hybrid}

The target route is not a pure Plan 3 single-level route.  It is:
\[
   \hbox{move inward along torsion levels}
   \quad\Longrightarrow\quad
   \hbox{discard an }O(\dT(\Om))\hbox{ layer}
   \quad\Longrightarrow\quad
   \hbox{enter the small }L^\infty\hbox{ radial-graph class}
\]
and then use the existing small-amplitude radial-graph theorem plus the
superlevel transfer.

Let \(u=u_\Om\), \(E_t=\{u>t\}\), and normalize \(|\Om|=\pi\) in
dimension two.  The exact theorem one would like is:

\begin{conjecture}[Planar \(O(\dT)\) graph-entry level]
\label{conj:entry}
There are computable dimensional constants \(\delta_0,\eta_0,C>0\) such
that if \(\Om\subset\R^2\) is open, \(|\Om|=\pi\), and
\(\dT(\Om)\le\delta_0\), then there is a regular level \(E_{\hat t}\) and a
center \(z\) such that
\[
   |\Om\setminus E_{\hat t}|\le C\dT(\Om),
\]
and, after translating by \(z\) and rescaling to area \(\pi\),
\[
   E_{\hat t}^{\#}
      =\{re^{i\theta}:0<r<1+\varphi(\theta)\},
   \qquad
   \|\varphi\|_{L^\infty(S^1)}\le\eta_0 .
\]
\end{conjecture}

If Conjecture~\ref{conj:entry} is true, the rest is already present in
the repo.  Step 3 gives
\[
   \dT(E_{\hat t})\le (1-C\dT(\Om))^{-1}\dT(\Om)\le 2\dT(\Om),
\]
the direct radial-graph theorem
\texttt{H12\_ASYMMETRY\_DIRECT\_PROOF.tex} gives
\[
   \A(E_{\hat t})^2\le C\dT(E_{\hat t}),
\]
and Step 5 gives
\[
   \A(\Om)\le \A(E_{\hat t})+
      2\frac{|\Om\setminus E_{\hat t}|}{|\Om|}.
\]
Thus \(\A(\Om)^2\le C\dT(\Om)\).  This closure uses no BDV theorem.

\section{What the repo supplies}

The following ingredients are usable.

\begin{enumerate}
\item \textbf{Good near-boundary level.}  The corrected Step 1
\texttt{proof-step1.md} works in the volume variable and gives a
regular level with
\[
   |\Om\setminus E_{\hat t}|\le C_\theta \dT(\Om),
   \qquad
   D_H(\hat t)\le C\theta ,
\]
for fixed absolute \(\theta>0\).  The same Markov selection can be made
for \(D_H+D_I\) at the same fixed scale.  The output is a small layer
and a small weighted Cauchy--Schwarz defect, not a pointwise gradient
floor or a graph.

\item \textbf{Exact superlevel deficit transfer.}  Step 3 is sound:
\(u-\hat t\) is exactly the torsion function of \(E_{\hat t}\), and the
profile computation gives the displayed bound on \(\dT(E_{\hat t})\).

\item \textbf{Transfer back.}  Step 5 is a purely measure-theoretic
triangle inequality for Fraenkel asymmetry and is sound.

\item \textbf{Small-amplitude radial graphs.}  The direct radial-graph
proof is the right final theorem: it requires a genuine one-crossing
radial graph and small \(L^\infty\) amplitude, but no slope bound.

\item \textbf{Plan 2 slicing technology.}  The Maggi radial slicing
identity and the bad-ray multiplicity absorption in
\texttt{wave2-A-radial-l2-trace.md} are useful.  In a fixed annular
bracket, the number of extra crossings on rays is linearly charged by
perimeter excess.  This is a real tool, but it controls the measure of
bad rays; it does not make the bad-ray set empty.

\item \textbf{Self-truncated velocity.}  The note
\texttt{WIP\_TrimmedVelocityRepair.tex} correctly extracts from
\(D_H\) that low-gradient portions of a good level are paid for by the
\(1/|\nabla u|\) weight.  This is the right mechanism for active thin
outward fingers in dimension two.
\end{enumerate}

The following cannot be used as load-bearing facts.

\begin{enumerate}
\item The old unweighted Hessian Chain B is false.  The high-frequency
near-ball computation in \texttt{PLAN3\_N2\_SINGLE\_LEVEL\_RECHECK.tex}
shows that no dimensional inequality
\[
   \dT(D)\ge c |D|^{-1}\int_D |D^2u_D+I/2|^2
\]
can hold on all smooth near-balls.

\item The Plan 2 endpoint metric trace is not a literal graph-entry
theorem.  Even where it gives \(d_{\mathcal X}(E_\rho,B_\rho)^2\le
C\dT\), this is \(L^1\)-modulo-translation control, not a one-crossing
\(L^\infty\) radial parametrization.

\item The WIP same-centre FJ package is useful for centering and normal
oscillation, but it is not a proof of Conjecture~\ref{conj:entry}.
It still tolerates small bad-ray sets and imports isoperimetric
near-spherical/selection facts, not torsion graph entry.
\end{enumerate}

\section{A better replacement for literal graph entry}

The correct repair target is slightly weaker than
Conjecture~\ref{conj:entry} and appears more realistic.

\begin{conjecture}[Planar \(O(\dT)\) radialization lemma]
\label{conj:radialize}
Let \(E=E_{\hat t}\) be the Step-1 regular level in dimension two.
There exists a set \(G\), after a translation and a rescaling, such that
\[
   G=\{re^{i\theta}:0<r<1+\varphi(\theta)\},
   \qquad
   \|\varphi\|_{L^\infty}\le\eta_0,
\]
and
\[
   |E\Delta G|\le C\dT(\Om),
   \qquad
   \dT(G)\le C\dT(\Om).
   \tag{R}
\]
\end{conjecture}

This is enough.  The graph theorem gives \(\A(G)^2\le C\dT(G)\le
C\dT(\Om)\), while
\[
   \A(E)\le \A(G)+C|E\Delta G|
\]
and Step 5 transfers from \(E\) to \(\Om\).  Notice the important
point: the repair volume in (R) must be \(O(\dT)\), not merely
\(O(\sqrt{\dT})\), unless one has a separate monotonicity argument
showing that the repair does not increase the deficit beyond
\(C\dT\).

Conjecture~\ref{conj:radialize} is the most promising form of the
missing piece because it allows filling tiny holes or deleting tiny
loops.  Literal graph entry for the actual level is too rigid: a tiny
central hole already destroys one-crossing radiality on every ray, even
though its volume may be much smaller than the torsion deficit.

\section{Attempt to prove the radialization lemma}

The natural decomposition is along rays from a same center \(z\), using
Plan 2's slicing notation.  For a.e. angle \(\theta\), write
\[
   E_\theta=\{r>0:z+r e^{i\theta}\in E\},
   \qquad
   \partial^*E_\theta=\{r_{\theta,1}<\cdots<r_{\theta,N(\theta)}\}.
\]
Good rays have \(N(\theta)=1\); bad rays have \(N(\theta)\ge2\).

\subsection{Good rays}

On good rays the radial graph candidate is clear: keep the first
crossing \(r_{\theta,1}\).  The Plan 2 direct \(T_1\) slicing bound
controls the angular \(L^1\) radial mismatch by the symmetric difference
and perimeter excess.  This is useful for metric/Fraenkel closure, but
it is not \(L^\infty\) control.  To obtain small amplitude one still
needs an annular containment theorem:
\[
   B_{1-\eta}(z)\subset E \subset B_{1+\eta}(z),
   \qquad \eta\ll1.
   \tag{A}
\]
For the Step-1 level, (A) is not proved in the repo.  FMP gives
measure closeness at the right square-root rate when a pointwise
isoperimetric deficit is known; it does not give Hausdorff annular
containment for arbitrary finite-perimeter levels.  In dimension two,
one expects fixed-depth inward slits to cost torsion through capacity
or perimeter, but that statement is not currently proved.

\subsection{Bad rays}

The Plan 2 bad-ray estimate is genuinely helpful:
\[
   \int_{S^1} (N(\theta)-1)\,d\theta
      \le C\,\bigl(P(E)-P(B)\bigr)
\]
inside an annular bracket.  Thus multiplicity is small in angular
measure when the isoperimetric deficit of the level is small.

This still falls short of Conjecture~\ref{conj:radialize}.  First, the
Step-1 \(O(\dT)\)-layer selection only gives \(D_I\) small at a fixed
absolute threshold if one keeps the layer \(O(\dT)\).  It does not give
\(D_I(E_{\hat t})=O(\dT)\).  Second, even an \(O(\dT)\) angular measure
bound would not by itself give a graph; it would only say that the
non-graph rays form a small set.  One must also repair those rays and
prove both parts of (R): small symmetric-difference cost and no loss in
the torsion deficit.

\subsection{Outward fingers}

Here dimension two is favorable.  On a thin outward tube of width \(w\)
and length \(L\), active tube levels have
\[
   D_H(t)\simeq \frac{L}{w}.
\]
Thus a Step-1 level with \(D_H\le\theta\) cannot carry a long thin
active finger unless \(L\lesssim \theta w\).  In particular, the volume
of such a finger is \(wL\lesssim \theta w^2\), and the extra layer
needed to clear the tube is of the same \(w^2\) order.  This is exactly
where \(n=2\) differs from higher dimensions.

This model computation is strong evidence that outward fingers can be
radialized at \(O(\dT)\) cost, but the repo does not contain the
general geometric lemma turning this model into a proof for arbitrary
bad components of an analytic torsion level.

\subsection{Inward fjords and holes}

Moving inward cannot fill a missing fjord or a hole.  Therefore a
literal superlevel graph-entry theorem must prove that no fixed-depth
inward defect remains.  The realistic statement is instead that the
defect can be filled at \(O(\dT)\) cost while preserving
\(\dT(G)\le C\dT(\Om)\).

The expected two-dimensional tool is capacity/modulus.  A slit or fjord
of fixed radial depth should force a definite torsion loss, uniformly in
its width, while a small hole of radius \(r\) has torsion cost of order
\(1/|\log r|\), much larger than its area \(r^2\).  This would imply
that the area to be filled is \(O(\dT)\), often much smaller.

But this is still a theorem to prove.  The current repo only formulates
this capacity idea; it does not provide the needed quantitative
decomposition for arbitrary same-ray cancellations.

\section{Outcome}

The remaining puzzle piece has not been proved by the existing notes.
The repo does contain the right components, but they stop one step
short of the needed theorem.

The cleanest route forward is:

\begin{enumerate}
\item Prove a planar annular-entry/capacity lemma: every fixed-depth
inward fjord or hole in a near-optimal torsion level costs a computable
positive amount of \(\dT\); quantitatively, the total volume that must
be filled to obtain annular containment is \(O(\dT)\).

\item Prove a planar bad-ray radialization lemma in a thin annulus:
using Maggi slicing and the Plan 2 multiplicity absorption, the total
volume of same-ray multi-crossing/cancellation that must be modified is
\(O(\dT)\), and the modification satisfies \(\dT(G)\le C\dT(E)\).

\item Feed the resulting radial graph \(G\) into
\texttt{H12\_ASYMMETRY\_DIRECT\_PROOF.tex} and transfer back.
\end{enumerate}

This is still promising, but it is not closed.  The underclaim in the
older notes was thinking that annular or metric \(L^1\) control is too
weak to be useful; it is useful if one allows \(O(\dT)\) radialization.
The overclaim would be saying that the repo already proves this
radialization.  It does not.  The missing theorem is exactly the
capacity-plus-bad-ray radialization lemma above.

\end{document}
